\documentclass[11pt]{article}
\usepackage[T1]{fontenc}
\usepackage{amsmath,amssymb}
\usepackage{cite}
\usepackage{booktabs}
\usepackage{graphicx}
\usepackage{xurl}
\usepackage{hyperref}
\hypersetup{
  colorlinks=true,
  linkcolor=blue,
  urlcolor=blue,
  citecolor=blue,
  pdfauthor={Jabulani Monchwe, Martin Chanza, Olorato Motitswane, Mfundo Monchwe}
}
\newcommand{\software}{\textsf{ReliableML}}
\newcommand{\code}[1]{\texttt{\small #1}}
\newcommand{\jabulaniauthor}{
\IEEEauthorblockN{Jabulani Monchwe}
\IEEEauthorblockA{\textit{Applied Business Analytics and Decision Making}\\
\textit{North-West University}\\
\textit{IDEA-Research Club Chairperson}\\
Johannesburg, South Africa\\
\href{mailto:monchwej.25@gmail.com}{monchwej.25@gmail.com}}
}
\newcommand{\martinchanzauthor}{
\IEEEauthorblockN{Martin Chanza}
\IEEEauthorblockA{\textit{Applied Business Analytics and Decision Making}\\
\textit{North-West University}\\
Mafikeng, South Africa\\
\href{mailto:martin.chanza@nwu.ac.za}{martin.chanza@nwu.ac.za}}
}
\newcommand{\oloratoauthor}{
\IEEEauthorblockN{Olorato Motitswane}
\IEEEauthorblockA{\textit{Eduvos Midrand Campus -- IT Faculty Data Science, South Africa}\\
\textit{IDEA-Research Club}\\
\href{mailto:oloratomotitswane@gmail.com}{oloratomotitswane@gmail.com}}
}
\newcommand{\mfundoauthor}{
\IEEEauthorblockN{Mfundo Monchwe}
\IEEEauthorblockA{\textit{Eduvos Midrand Campus -- IT Faculty Data Science, South Africa}\\
\textit{IDEA-Research Club}\\
\href{mailto:diditmfundo@gmail.com}{diditmfundo@gmail.com}}
}
\newcommand{\minenhleauthor}{
\IEEEauthorblockN{Minenhle Mchunu}
\IEEEauthorblockA{\textit{Applied Business Analytics and Decision Making}\\
\textit{Eduvos Midrand Campus -- IT Faculty Data Science, South Africa}\\
\textit{IDEA-Research Club}\\
\href{mailto:Mchunu.minenhle02@gmail.com}{Mchunu.minenhle02@gmail.com}}
}
\newcommand{\projecturl}{\url{https://github.com/Jabualani/ML-Reliabilty}}

\title{\software{} Method Note: Statistical Reason Audit}
\author{\contactauthors}
\date{Version 0.1.1}
\begin{document}
\maketitle
\begin{abstract}
The reason audit combines labeled validity, perturbation-based feature
reliance, subgroup and time stability, counterfactual sensitivity, logistic
associations, and explanation alignment. It investigates potential
shortcuts; it does not prove causality or correct reasons.
\end{abstract}

\section{Audit design}
Use labeled, representative holdout rows not used for fitting or model
selection, plus training/reference background data. The report includes
bootstrap intervals for predictive measures such as accuracy, multiclass
Brier score, and top-label calibration error. This addresses whether
predictions perform adequately on the evaluated sample, independently of
what explanation methods say.

\section{Feature reliance}
For each feature, the audit perturbs held-out values using donor rows and
measures the increase in held-out Brier loss. Conditional donor selection can
use nearest neighbors over user-specified numeric covariates; this is a
conditional-distribution approximation and can fail in sparse regions.
Without conditional covariates, marginal perturbation may construct
unrealistic combinations when features are correlated. Bootstrap intervals
and repeated estimates quantify sampling variability but not all sources of
uncertainty.

\section{Stability and sensitivity}
Optional groups and time periods report reliance across slices. Optional
counterfactual replacements report prediction-flip rates and probability
changes. The user is responsible for supplying domain-valid changes; arbitrary
edits do not represent plausible interventions. Negative-control features
and features unavailable at prediction time produce warnings, not definitive
proof of leakage or shortcut learning.

\section{Association and explanation alignment}
For binary or selected one-vs-rest outcomes, an unregularized logistic model
reports standardized-feature odds ratios, confidence intervals, p-values, and
Benjamini--Hochberg adjusted p-values. These quantify conditional associations
with observed labels. The audit can compare row-level signed Shapley
contributions with logistic log-odds contributions and compare supplied
importance mappings with global Shapley ranks. Agreement is not a test of
causality and can occur under confounding or shared bias.

\section{Interpretation}
No single score in the audit certifies correct-for-correct-reasons prediction.
Conclusions depend on holdout representativeness, label quality, measurement,
feature availability, dependence structure, calibration, and causal/domain
knowledge. Current bootstrap procedures treat rows as independent; grouped
or time-dependent data may need block or cluster resampling. Audits are
computationally intensive because they repeatedly perturb data and evaluate
the model.

\section{Interface}
\begin{verbatim}
audit = model.audit_reasons(
    X_holdout, y_holdout, positive_label=1,
    conditional_on=["age", "region"],
    groups={"sex": sex_holdout},
    time_periods=month_holdout,
    counterfactuals={"income": plausible_income_changes},
)
print(audit.predictive_validity)
print(audit.feature_reliance)
print(audit.logistic_association)
\end{verbatim}

\section{Further reading}
\begin{thebibliography}{3}
\bibitem{shap}
S. M. Lundberg and S.-I. Lee. A unified approach to interpreting model
predictions. \emph{Advances in Neural Information Processing Systems}, 2017.
\url{https://arxiv.org/abs/1705.07874}.
\bibitem{hardt}
M. Hardt, E. Price, and N. Srebro. Equality of opportunity in supervised
learning. \emph{Advances in Neural Information Processing Systems}, 2016.
\url{https://arxiv.org/abs/1610.02413}.
\bibitem{reliability}
\software{} reason audit assumptions and interfaces:
\url{https://github.com/Jabualani/ML-Reliabilty}.
\end{thebibliography}
\end{document}
